\chapter{Complexity — Machine}
%%% generated by the blueprint pipeline from TCSlib.Complexity.ClassPSPACE.TQBFARM.Machine
\section{Overview}

Declarations of the Lean module \texttt{TCSlib.Complexity.ClassPSPACE.TQBFARM.Machine}.

\section{Declarations}

\begin{definition}[Pair of input languages]
\lean{Complexity.LogProg.TQBFARM.decs}
\label{Complexity.LogProg.TQBFARM.decs}
\leanok
\uses{Complexity.inputBitLang, Complexity.inputLenLang}
Define a function from the two-element index set to languages over Boolean lists by
assigning the language of input bits to index $0$ and the language of input length
encodings to index $1$.
\end{definition}

\begin{lemma}[The Two Index Languages Are Logarithmic-Space Decidable]
\lean{Complexity.LogProg.TQBFARM.decs_mem}
\label{Complexity.LogProg.TQBFARM.decs_mem}
\leanok
\uses{Complexity.LOGSPACE, Complexity.LogProg.TQBFARM.decs, Complexity.inputBitLang_mem, Complexity.inputLenLang_mem}
\difficulty{4}
Each of the two languages that, on an encoded pair $\langle u,\mathrm{bits}(i)\rangle$,
respectively tests whether the $i$th bit of $u$ is true (with a missing bit treated as
false) or whether $i<|u|$, is decidable in logarithmic space.
\end{lemma}

\begin{definition}[Language membership indicators]
\lean{Complexity.LogProg.TQBFARM.orc}
\label{Complexity.LogProg.TQBFARM.orc}
\leanok
\uses{Complexity.LogProg.TQBFARM.decs, Turing.MultiTapeTM.indicator}
For each index $j$ in the two-element set and each Boolean list $V$, define the Boolean
indicator of whether $V$ belongs to the language assigned to $j$.
\end{definition}

\begin{lemma}[Oracle query for an input bit]
\lean{Complexity.LogProg.TQBFARM.orc_bit}
\label{Complexity.LogProg.TQBFARM.orc_bit}
\leanok
\uses{Complexity.LogProg.TQBFARM.orc, Complexity.indicator_inputBitLang, Complexity.inputBitLang, Turing.MultiTapeTM.indicator, Turing.pairEncode}
\difficulty{3}
For every bit string $y$ and natural number $a$, the first oracle returns the bit of $y$
at position $a$ on the encoded query $\langle y,\mathrm{bits}(a)\rangle$, taking the
value false when $a$ is outside the string.
\end{lemma}

\begin{lemma}[Oracle query for input length]
\lean{Complexity.LogProg.TQBFARM.orc_len}
\label{Complexity.LogProg.TQBFARM.orc_len}
\leanok
\uses{Complexity.LogProg.TQBFARM.orc, Complexity.indicator_inputLenLang, Complexity.inputLenLang, Turing.MultiTapeTM.indicator, Turing.pairEncode}
\difficulty{3}
For every bit string $y$ and natural number $a$, the second oracle returns true on
$\langle y,\mathrm{bits}(a)\rangle$ exactly when $a<|y|$, and false otherwise.
\end{lemma}

\begin{definition}[First element of a seven-element set]
\lean{Complexity.LogProg.TQBFARM.rK}
\label{Complexity.LogProg.TQBFARM.rK}
\leanok
Define $r_K$ to be the element $0$ of the set of integers from $0$ through $6$.
\end{definition}

\begin{definition}[Second element of a seven-element set]
\lean{Complexity.LogProg.TQBFARM.rS}
\label{Complexity.LogProg.TQBFARM.rS}
\leanok
Define $r_S$ to be the element $1$ of the set of integers from $0$ through $6$.
\end{definition}

\begin{definition}[Third element of a seven-element set]
\lean{Complexity.LogProg.TQBFARM.rP}
\label{Complexity.LogProg.TQBFARM.rP}
\leanok
Define $r_P$ to be the element $2$ of the set of integers from $0$ through $6$.
\end{definition}

\begin{definition}[Fourth element of a seven-element set]
\lean{Complexity.LogProg.TQBFARM.rT}
\label{Complexity.LogProg.TQBFARM.rT}
\leanok
Define $r_T$ to be the element $3$ of the set of integers from $0$ through $6$.
\end{definition}

\begin{definition}[Fifth element of a seven-element set]
\lean{Complexity.LogProg.TQBFARM.rZ}
\label{Complexity.LogProg.TQBFARM.rZ}
\leanok
Define $r_Z$ to be the element $4$ of the set of integers from $0$ through $6$.
\end{definition}

\begin{definition}[Sixth element of a seven-element set]
\lean{Complexity.LogProg.TQBFARM.rC}
\label{Complexity.LogProg.TQBFARM.rC}
\leanok
Define $r_C$ to be the element $5$ of the set of integers from $0$ through $6$.
\end{definition}

\begin{definition}[Variable register index]
\lean{Complexity.LogProg.TQBFARM.rV}
\label{Complexity.LogProg.TQBFARM.rV}
\leanok
Let the seven register positions be indexed by $0,1,\ldots,6$. The variable register is
assigned index $6$.
\end{definition}

\begin{definition}[Machine state]
\lean{Complexity.LogProg.TQBFARM.St}
\label{Complexity.LogProg.TQBFARM.St}
\leanok
A machine state consists of seven natural numbers: a depth, an assignment stack, an
input pointer, a scratch copy of the stack, a scratch register for copy loops, a halving
counter, and the variable of the current literal.
\end{definition}

\begin{definition}[State register vector]
\lean{Complexity.LogProg.TQBFARM.St.toV}
\label{Complexity.LogProg.TQBFARM.St.toV}
\leanok
\uses{Complexity.LogProg.TQBFARM.St}
For a machine state, form a vector of seven natural numbers whose entries, in order, are
its depth, assignment stack, input pointer, scratch copy, copy-loop scratch register,
halving counter, and current literal's variable.
\end{definition}

\begin{definition}[Zero state]
\lean{Complexity.LogProg.TQBFARM.St.zero}
\label{Complexity.LogProg.TQBFARM.St.zero}
\leanok
\uses{Complexity.LogProg.TQBFARM.St}
The zero state is the machine state in which all seven registers are zero.
\end{definition}

\begin{lemma}[Zero state register vector]
\lean{Complexity.LogProg.TQBFARM.St.toV_zero}
\label{Complexity.LogProg.TQBFARM.St.toV_zero}
\leanok
\uses{Complexity.LogProg.TQBFARM.St.toV, Complexity.LogProg.TQBFARM.St.zero}
\difficulty{4}
The register vector of the zero state is the constant function on the seven registers
with value $0$.
\end{lemma}

\begin{lemma}[Depth Coordinate]
\lean{Complexity.LogProg.TQBFARM.St.toV_K}
\label{Complexity.LogProg.TQBFARM.St.toV_K}
\leanok
\uses{Complexity.LogProg.TQBFARM.St, Complexity.LogProg.TQBFARM.St.toV, Complexity.LogProg.TQBFARM.rK}
\difficulty{1}
For every state, the coordinate of its register vector indexed by the depth register is
equal to the value of its depth register.
\end{lemma}

\begin{lemma}[Assignment Stack Coordinate]
\lean{Complexity.LogProg.TQBFARM.St.toV_S}
\label{Complexity.LogProg.TQBFARM.St.toV_S}
\leanok
\uses{Complexity.LogProg.TQBFARM.St, Complexity.LogProg.TQBFARM.St.toV, Complexity.LogProg.TQBFARM.rS}
\difficulty{1}
For every state, the coordinate of its register vector indexed by the assignment stack
is equal to the value of its assignment stack.
\end{lemma}

\begin{lemma}[Input Pointer Coordinate]
\lean{Complexity.LogProg.TQBFARM.St.toV_P}
\label{Complexity.LogProg.TQBFARM.St.toV_P}
\leanok
\uses{Complexity.LogProg.TQBFARM.St, Complexity.LogProg.TQBFARM.St.toV, Complexity.LogProg.TQBFARM.rP}
\difficulty{1}
For every state, the coordinate of its register vector indexed by the input pointer is
equal to the value of its input pointer.
\end{lemma}

\begin{lemma}[Scratch Copy Coordinate]
\lean{Complexity.LogProg.TQBFARM.St.toV_T}
\label{Complexity.LogProg.TQBFARM.St.toV_T}
\leanok
\uses{Complexity.LogProg.TQBFARM.St, Complexity.LogProg.TQBFARM.St.toV, Complexity.LogProg.TQBFARM.rT}
\difficulty{1}
For every state, the coordinate of its register vector indexed by the scratch copy
register is equal to the value of that register.
\end{lemma}

\begin{lemma}[Scratch register lookup]
\lean{Complexity.LogProg.TQBFARM.St.toV_Z}
\label{Complexity.LogProg.TQBFARM.St.toV_Z}
\leanok
\uses{Complexity.LogProg.TQBFARM.St, Complexity.LogProg.TQBFARM.St.toV, Complexity.LogProg.TQBFARM.rZ}
\difficulty{1}
Let a state consist of seven natural-number registers: depth, assignment stack, input
pointer, scratch copy, copy-loop scratch register, halving counter, and current-literal
variable. In the register vector ordered by these registers, the entry at index $4$
equals the value of the scratch register.
\end{lemma}

\begin{lemma}[Scratch Counter Coordinate]
\lean{Complexity.LogProg.TQBFARM.St.toV_C}
\label{Complexity.LogProg.TQBFARM.St.toV_C}
\leanok
\uses{Complexity.LogProg.TQBFARM.St, Complexity.LogProg.TQBFARM.St.toV, Complexity.LogProg.TQBFARM.rC}
\difficulty{1}
For every state, the coordinate of its register vector indexed by the counter register
is equal to the value of its counter register.
\end{lemma}

\begin{lemma}[Current Variable Coordinate]
\lean{Complexity.LogProg.TQBFARM.St.toV_V}
\label{Complexity.LogProg.TQBFARM.St.toV_V}
\leanok
\uses{Complexity.LogProg.TQBFARM.St, Complexity.LogProg.TQBFARM.St.toV, Complexity.LogProg.TQBFARM.rV}
\difficulty{1}
For every state, the coordinate of its register vector indexed by the current-variable
register is equal to the value of that register.
\end{lemma}

\begin{lemma}[Updating the depth register]
\lean{Complexity.LogProg.TQBFARM.St.upd_K}
\label{Complexity.LogProg.TQBFARM.St.upd_K}
\leanok
\uses{Complexity.LogProg.TQBFARM.St, Complexity.LogProg.TQBFARM.St.toV, Complexity.LogProg.TQBFARM.rK}
\difficulty{4}
For any machine state $s$ and natural number $k$, updating register $0$ of the register
vector of $s$ to $k$ gives the register vector of the state obtained from $s$ by setting
its depth $K$ to $k$.
\end{lemma}

\begin{lemma}[Updating the assignment stack]
\lean{Complexity.LogProg.TQBFARM.St.upd_S}
\label{Complexity.LogProg.TQBFARM.St.upd_S}
\leanok
\uses{Complexity.LogProg.TQBFARM.St, Complexity.LogProg.TQBFARM.St.toV, Complexity.LogProg.TQBFARM.rS}
\difficulty{4}
For any machine state $s$ and natural number $k$, updating register $1$ of the register
vector of $s$ to $k$ gives the register vector of the state obtained from $s$ by setting
its assignment stack $S$ to $k$.
\end{lemma}

\begin{lemma}[Updating the input pointer]
\lean{Complexity.LogProg.TQBFARM.St.upd_P}
\label{Complexity.LogProg.TQBFARM.St.upd_P}
\leanok
\uses{Complexity.LogProg.TQBFARM.St, Complexity.LogProg.TQBFARM.St.toV, Complexity.LogProg.TQBFARM.rP}
\difficulty{4}
For any machine state $s$ and natural number $k$, updating register $2$ of the register
vector of $s$ to $k$ gives the register vector of the state obtained from $s$ by setting
its input pointer $P$ to $k$.
\end{lemma}

\begin{lemma}[Updating the scratch stack copy]
\lean{Complexity.LogProg.TQBFARM.St.upd_T}
\label{Complexity.LogProg.TQBFARM.St.upd_T}
\leanok
\uses{Complexity.LogProg.TQBFARM.St, Complexity.LogProg.TQBFARM.St.toV, Complexity.LogProg.TQBFARM.rT}
\difficulty{4}
For any machine state $s$ and natural number $k$, updating register $3$ of the register
vector of $s$ to $k$ gives the register vector of the state obtained from $s$ by setting
its scratch copy of the stack $T$ to $k$.
\end{lemma}

\begin{lemma}[Updating the scratch register]
\lean{Complexity.LogProg.TQBFARM.St.upd_Z}
\label{Complexity.LogProg.TQBFARM.St.upd_Z}
\leanok
\uses{Complexity.LogProg.TQBFARM.St, Complexity.LogProg.TQBFARM.St.toV, Complexity.LogProg.TQBFARM.rZ}
\difficulty{4}
For every machine state $st$ and natural number $k$, changing coordinate $4$ of the
register vector associated with $st$ to $k$ gives the register vector associated with
the state obtained by replacing its scratch register $Z$ by $k$.
\end{lemma}

\begin{lemma}[Updating the halving counter]
\lean{Complexity.LogProg.TQBFARM.St.upd_C}
\label{Complexity.LogProg.TQBFARM.St.upd_C}
\leanok
\uses{Complexity.LogProg.TQBFARM.St, Complexity.LogProg.TQBFARM.St.toV, Complexity.LogProg.TQBFARM.rC}
\difficulty{4}
For any machine state $s$ and natural number $k$, updating register $5$ of the register
vector of $s$ to $k$ gives the register vector of the state obtained from $s$ by setting
its halving counter $C$ to $k$.
\end{lemma}

\begin{lemma}[Updating the current literal variable]
\lean{Complexity.LogProg.TQBFARM.St.upd_V}
\label{Complexity.LogProg.TQBFARM.St.upd_V}
\leanok
\uses{Complexity.LogProg.TQBFARM.St, Complexity.LogProg.TQBFARM.St.toV, Complexity.LogProg.TQBFARM.rV}
\difficulty{4}
For any machine state $s$ and natural number $k$, updating register $6$ of the register
vector of $s$ to $k$ gives the register vector of the state obtained from $s$ by setting
the variable of the current literal $V$ to $k$.
\end{lemma}

\begin{definition}[Bounded machine state]
\lean{Complexity.LogProg.TQBFARM.St.Bd}
\label{Complexity.LogProg.TQBFARM.St.Bd}
\leanok
\uses{Complexity.LogProg.TQBFARM.St}
For a natural number $B$ and a machine state, the state is bounded by $B$ when each of
its seven registers is at most $B$.
\end{definition}

\begin{lemma}[Register Bounds Under the State Encoding]
\lean{Complexity.LogProg.TQBFARM.St.Bd.toV}
\label{Complexity.LogProg.TQBFARM.St.Bd.toV}
\leanok
\uses{Complexity.LogProg.TQBFARM.St, Complexity.LogProg.TQBFARM.St.Bd, Complexity.LogProg.TQBFARM.St.toV}
\difficulty{4}
For any natural number $B$ and state $\mathrm{st}$, if each of its seven registers is at
most $B$, then every coordinate of its register vector is at most $B$.
\end{lemma}

\begin{definition}[Copy-loop sites]
\lean{Complexity.LogProg.TQBFARM.CpS}
\label{Complexity.LogProg.TQBFARM.CpS}
\leanok
The copy-loop site set has six kinds of labels: two node-processing sites, two return
sites carrying a Boolean child value, and two literal-processing sites carrying a
Boolean polarity. The first node site copies the depth into the input pointer; the
second does the same for the second copy, for which the pointer is $2K$. The return
sites make the corresponding first and second copies, and the literal sites copy the
depth into the counter or the stack into the scratch register.
\end{definition}

\begin{definition}[Machine control labels]
\lean{Complexity.LogProg.TQBFARM.Lb}
\label{Complexity.LogProg.TQBFARM.Lb}
\leanok
\uses{Complexity.LogProg.TQBFARM.CpS, Complexity.LogProg.TQBFARM.SynSt}
The control-label set consists of an initialization label; syntax-checking labels
parameterized by automaton states and input bits; node, push, return, quantifier, leaf,
matrix, clause, and literal-processing labels with the Boolean parameters indicated by
their roles; and copy, doubling, move, subtraction, and halving loop labels. The copy
labels carry a copy-loop site and a phase in $\{0,\ldots,6\}$; doubling labels carry a
Boolean and a phase in $\{0,\ldots,3\}$; move, subtraction, and halving labels each
carry a Boolean and a phase in $\{0,1,2\}$.
\end{definition}

\begin{definition}[Copy source register]
\lean{Complexity.LogProg.TQBFARM.cpSrc}
\label{Complexity.LogProg.TQBFARM.cpSrc}
\leanok
\uses{Complexity.LogProg.TQBFARM.CpS, Complexity.LogProg.TQBFARM.rK, Complexity.LogProg.TQBFARM.rS}
For each copy-loop site, the source register is the depth register (index $0$), except
at the literal site that copies the stack, where the source register is the assignment
stack (index $1$).
\end{definition}

\begin{definition}[Copy destination register]
\lean{Complexity.LogProg.TQBFARM.cpDst}
\label{Complexity.LogProg.TQBFARM.cpDst}
\leanok
\uses{Complexity.LogProg.TQBFARM.CpS, Complexity.LogProg.TQBFARM.rC, Complexity.LogProg.TQBFARM.rP, Complexity.LogProg.TQBFARM.rT}
The destination register for each copy-loop site is $P$ at the two node-copy sites and
both return-copy sites, $C$ when copying the depth for a literal, and $T$ when copying
the stack for a literal.
\end{definition}

\begin{definition}[Copy-loop exit states]
\lean{Complexity.LogProg.TQBFARM.cpExit}
\label{Complexity.LogProg.TQBFARM.cpExit}
\leanok
\uses{Complexity.LogProg.TQBFARM.CpS, Complexity.LogProg.TQBFARM.Lb}
The exit state after the copy loop is the phase-zero loop at the second node-copy site
after the first node-copy site, and the node-reading state after the second node-copy
site. At the return-copy sites it is the phase-zero loop at the second return-copy site,
carrying the same child value, then the quantifier-reading state carrying that value;
after copying depth for a literal it is the phase-zero subtraction loop carrying the
polarity, and after copying the stack it is the phase-zero halving loop carrying the
polarity.
\end{definition}

\begin{definition}[Move-loop instruction]
\lean{Complexity.LogProg.TQBFARM.mvI}
\label{Complexity.LogProg.TQBFARM.mvI}
\leanok
\uses{Complexity.LogProg.Ins, Complexity.LogProg.TQBFARM.Lb}
Given source and destination registers, an exit state, a function assigning a state to
each of three phases, and a phase $j$, the instruction is: test whether the source is
zero, branching to the exit state if so and otherwise to phase $1$ when $j=0$; decrement
the source and go to phase $2$ when $j=1$; or increment the destination and go to phase
$0$ when $j=2$.
\end{definition}

\begin{definition}[Doubling-loop instruction]
\lean{Complexity.LogProg.TQBFARM.dmI}
\label{Complexity.LogProg.TQBFARM.dmI}
\leanok
\uses{Complexity.LogProg.Ins, Complexity.LogProg.TQBFARM.Lb}
Given source and destination registers, an exit state, a function assigning a state to
each of four phases, and a phase $j$, the instruction is: test whether the source is
zero, branching to the exit state if so and otherwise to phase $1$ when $j=0$; decrement
the source and go to phase $2$ when $j=1$; increment the destination and go to phase $3$
when $j=2$; or increment the destination and go to phase $0$ when $j=3$.
\end{definition}

\begin{definition}[Subtraction-loop instruction]
\lean{Complexity.LogProg.TQBFARM.sbI}
\label{Complexity.LogProg.TQBFARM.sbI}
\leanok
\uses{Complexity.LogProg.Ins, Complexity.LogProg.TQBFARM.Lb}
Given registers $q$ and $r$, an exit state, a function assigning a state to each of
three phases, and a phase $j$, the instruction is: test whether $q$ is zero, branching
to the exit state if so and otherwise to phase $1$ when $j=0$; decrement $q$ and go to
phase $2$ when $j=1$; or decrement $r$ and go to phase $0$ when $j=2$.
\end{definition}

\begin{definition}[Halving-loop instruction]
\lean{Complexity.LogProg.TQBFARM.hvI}
\label{Complexity.LogProg.TQBFARM.hvI}
\leanok
\uses{Complexity.LogProg.Ins, Complexity.LogProg.TQBFARM.Lb}
Given registers $c$ and $r$, an exit state, a function assigning a state to each of
three phases, and a phase $j$, the instruction is: test whether $c$ is zero, branching
to the exit state if so and otherwise to phase $1$ when $j=0$; decrement $c$ and go to
phase $2$ when $j=1$; or halve $r$ and go to phase $0$ when $j=2$.
\end{definition}

\begin{definition}[Register-copy loop instruction]
\lean{Complexity.LogProg.TQBFARM.cpI}
\label{Complexity.LogProg.TQBFARM.cpI}
\leanok
\uses{Complexity.LogProg.Ins, Complexity.LogProg.TQBFARM.Lb, Complexity.LogProg.TQBFARM.rZ}
Given seven registers, a designated source and destination register, an exit label, and
a label for each of seven phases, this definition assigns an instruction to each phase.
The instruction tests whether the source is zero, transfers one unit from the source to
the destination while it is nonzero, and then restores the source while consuming the
same number of units from the zero register; it exits when the zero register reaches
zero.
\end{definition}

\begin{definition}[Quantified Boolean formula machine]
\lean{Complexity.LogProg.TQBFARM.A}
\label{Complexity.LogProg.TQBFARM.A}
\leanok
\uses{Complexity.LogProg.ARM, Complexity.LogProg.TQBFARM.Lb, Complexity.LogProg.TQBFARM.cpDst, Complexity.LogProg.TQBFARM.cpExit, Complexity.LogProg.TQBFARM.cpI, Complexity.LogProg.TQBFARM.cpSrc, Complexity.LogProg.TQBFARM.dmI, Complexity.LogProg.TQBFARM.hvI, Complexity.LogProg.TQBFARM.mvI, Complexity.LogProg.TQBFARM.rC, Complexity.LogProg.TQBFARM.rK, Complexity.LogProg.TQBFARM.rP, Complexity.LogProg.TQBFARM.rS, Complexity.LogProg.TQBFARM.rT, Complexity.LogProg.TQBFARM.rV, Complexity.LogProg.TQBFARM.sbI, Complexity.LogProg.TQBFARM.synStep}
The function assigns a seven-register, two-decider instruction to each machine state. It
specifies the full control flow for initializing the stack, validating and traversing
the quantified formula, evaluating its matrix and literals, returning child values, and
carrying out the copy, doubling, move, subtraction, and halving loops; each branch,
register update, decider call, and Boolean return is given by the corresponding state
case.
\end{definition}

\begin{lemma}[Instruction precondition holds universally]
\lean{Complexity.LogProg.TQBFARM.preS_A}
\label{Complexity.LogProg.TQBFARM.preS_A}
\leanok
\uses{Complexity.LogProg.AConf, Complexity.LogProg.PreS, Complexity.LogProg.TQBFARM.A, Complexity.LogProg.TQBFARM.Lb, Complexity.LogProg.TQBFARM.cpI, Complexity.LogProg.TQBFARM.dmI, Complexity.LogProg.TQBFARM.hvI, Complexity.LogProg.TQBFARM.mvI, Complexity.LogProg.TQBFARM.sbI}
\difficulty{5}
For every bit string $x$ and every configuration of the seven-register program, the
selected instruction satisfies its precondition: register comparisons use distinct
registers, input-word comparisons are made only on inputs with the required valid
encoding, and oracle calls use pairwise distinct argument registers.
\end{lemma}

\begin{definition}[Bounded machine run]
\lean{Complexity.LogProg.TQBFARM.Run}
\label{Complexity.LogProg.TQBFARM.Run}
\leanok
\uses{Complexity.LogProg.AReach, Complexity.LogProg.GuardLe, Complexity.LogProg.TQBFARM.A, Complexity.LogProg.TQBFARM.Lb, Complexity.LogProg.TQBFARM.St, Complexity.LogProg.TQBFARM.St.toV, Complexity.LogProg.TQBFARM.orc}
For an input bit string $x$, a bound $B\in\mathbb N$, labels $l,l'$, and machine states
$s,s'$, a run from $(l,s)$ to $(l',s')$ means that the machine has a guarded transition
sequence from the configuration with label $l$ and state $s$ to the configuration with
label $l'$ and state $s'$, on input $x$, within bound $B$.
\end{definition}

\begin{lemma}[Register Bounds Imply the Configuration Guard]
\lean{Complexity.LogProg.TQBFARM.gb_of}
\label{Complexity.LogProg.TQBFARM.gb_of}
\leanok
\uses{Complexity.LogProg.GuardLe, Complexity.LogProg.TQBFARM.A, Complexity.LogProg.TQBFARM.Lb, Complexity.LogProg.TQBFARM.St, Complexity.LogProg.TQBFARM.St.Bd, Complexity.LogProg.TQBFARM.St.Bd.toV, Complexity.LogProg.TQBFARM.St.toV, Complexity.LogProg.TQBFARM.preS_A}
\difficulty{1}
Let $x$ be a Boolean input word, $B$ a natural number, $	ext{st}$ a machine state whose
seven registers are all at most $B$, and $l$ an instruction label. The configuration
with label $l$, register values given by $	ext{st}$, and no answer satisfies the machine
guard: any register-equality instruction compares distinct registers; input-comparison
instructions are used only on inputs of the required valid plain or paired form; and an
oracle call has no repeated register arguments. Its register values are all at most $B$.
\end{lemma}

\begin{lemma}[Prepending One Machine Step]
\lean{Complexity.LogProg.TQBFARM.step_of}
\label{Complexity.LogProg.TQBFARM.step_of}
\leanok
\uses{Complexity.LogProg.AConf, Complexity.LogProg.AReach, Complexity.LogProg.AReach.step, Complexity.LogProg.GuardLe, Complexity.LogProg.TQBFARM.A, Complexity.LogProg.TQBFARM.Lb, Complexity.LogProg.TQBFARM.St, Complexity.LogProg.TQBFARM.St.Bd, Complexity.LogProg.TQBFARM.St.toV, Complexity.LogProg.TQBFARM.gb_of, Complexity.LogProg.TQBFARM.orc, Complexity.LogProg.astep}
\difficulty{1}
Let $x$ be a Boolean word, let $B$ be a natural-number bound, let $l$ be a program
label, let $st$ be a seven-register state, and let $c$ be an abstract configuration. If
every register of $st$ is at most $B$ and there is a finite run from the configuration
obtained by taking one step of the program from $(l,st)$ to $c$, with every
configuration before the endpoint satisfying the program guard and register bound $B$,
then there is such a run from $(l,st)$ to $c$.
\end{lemma}

\begin{lemma}[Reachability from Equal Register States]
\lean{Complexity.LogProg.TQBFARM.fin_of}
\label{Complexity.LogProg.TQBFARM.fin_of}
\leanok
\uses{Complexity.LogProg.AReach, Complexity.LogProg.AReach.refl, Complexity.LogProg.GuardLe, Complexity.LogProg.TQBFARM.A, Complexity.LogProg.TQBFARM.Lb, Complexity.LogProg.TQBFARM.St, Complexity.LogProg.TQBFARM.St.toV, Complexity.LogProg.TQBFARM.orc}
\difficulty{1}
Let $x$ be a Boolean input word, $B$ a natural-number bound, $l$ an instruction label,
and $	ext{st},	ext{st}'$ machine states with $	ext{st}=	ext{st}'$. Starting at label $l$
with register values given by $	ext{st}$ and no answer, the machine can reach label $l$
with register values given by $	ext{st}'$ and no answer, along a finite run in which
every configuration before the endpoint satisfies the machine guard and has every
register at most $B$.
\end{lemma}

\begin{lemma}[Distinct registers in copy operations]
\lean{Complexity.LogProg.TQBFARM.cp_ne}
\label{Complexity.LogProg.TQBFARM.cp_ne}
\leanok
\uses{Complexity.LogProg.TQBFARM.CpS, Complexity.LogProg.TQBFARM.cpDst, Complexity.LogProg.TQBFARM.cpSrc, Complexity.LogProg.TQBFARM.rZ}
\difficulty{3}
For every copy-operation site $k$, its selected source and destination register indices
are distinct, and neither index is the scratch register used by the copy loops.
\end{lemma}

\begin{theorem}[Copying one register into another]
\lean{Complexity.LogProg.TQBFARM.cp_run}
\label{Complexity.LogProg.TQBFARM.cp_run}
\leanok
\uses{Complexity.LogProg.AReach, Complexity.LogProg.GuardLe, Complexity.LogProg.TQBFARM.A, Complexity.LogProg.TQBFARM.CpS, Complexity.LogProg.TQBFARM.St, Complexity.LogProg.TQBFARM.St.Bd, Complexity.LogProg.TQBFARM.St.Bd.toV, Complexity.LogProg.TQBFARM.St.toV, Complexity.LogProg.TQBFARM.cpDst, Complexity.LogProg.TQBFARM.cpExit, Complexity.LogProg.TQBFARM.cpSrc, Complexity.LogProg.TQBFARM.cp_ne, Complexity.LogProg.TQBFARM.orc, Complexity.LogProg.copy_run}
\difficulty{1}
Let $x$ be a finite Boolean word, $B$ a natural number, $k$ one of the six copy-loop
sites, and $st$ a machine state whose seven register values are all at most $B$. Suppose
the scratch register $Z$ is zero and the sum of the initial values in the selected
destination and source registers is at most $B$. Then the machine reaches the exit label
for $k$ from phase $0$, with the destination register equal to its initial value plus
the source register's initial value and all other registers unchanged. Every
configuration before the endpoint satisfies the machine's validity conditions and has
all register values at most $B$.
\end{theorem}

\begin{theorem}[Doubling a register into scratch space]
\lean{Complexity.LogProg.TQBFARM.dm_run}
\label{Complexity.LogProg.TQBFARM.dm_run}
\leanok
\uses{Complexity.LogProg.TQBFARM.A, Complexity.LogProg.TQBFARM.Run, Complexity.LogProg.TQBFARM.St, Complexity.LogProg.TQBFARM.St.Bd, Complexity.LogProg.TQBFARM.St.Bd.toV, Complexity.LogProg.TQBFARM.St.toV, Complexity.LogProg.TQBFARM.orc, Complexity.LogProg.TQBFARM.rS, Complexity.LogProg.TQBFARM.rT, Complexity.LogProg.dmove_run}
\difficulty{3}
Let $x$ be a finite Boolean word, $B$ a natural number, $b$ a Boolean value, and $st$ a
machine state whose seven register values are all at most $B$. If $st.T+2st.S\le B$,
then the machine reaches label $\mathrm{mv}(b,0)$ from label $\mathrm{dm}(b,0)$ with
state obtained by setting $S$ to zero and replacing $T$ by its initial value plus twice
the initial value of $S$, leaving the other state coordinates unchanged. Every
configuration before the endpoint satisfies the machine's validity conditions and has
all register values at most $B$.
\end{theorem}

\begin{theorem}[Run of the register transfer loop]
\lean{Complexity.LogProg.TQBFARM.mv_run}
\label{Complexity.LogProg.TQBFARM.mv_run}
\leanok
\uses{Complexity.LogProg.TQBFARM.A, Complexity.LogProg.TQBFARM.Run, Complexity.LogProg.TQBFARM.St, Complexity.LogProg.TQBFARM.St.Bd, Complexity.LogProg.TQBFARM.St.Bd.toV, Complexity.LogProg.TQBFARM.St.toV, Complexity.LogProg.TQBFARM.orc, Complexity.LogProg.TQBFARM.rS, Complexity.LogProg.TQBFARM.rT, Complexity.LogProg.move_run}
\difficulty{3}
Let $x$ be a bit string, $B\in\mathbb{N}$, $b\in\{0,1\}$, and let $st=(K,S,P,T,Z,C,V)$
be a seven-register state. If every register is at most $B$ and $S+T\leq B$, the machine
has a guarded run from label $mv(b,0)$ with state $st$ to label $incS$ when $b$ is true
and $incK$ when $b$ is false, with state $(K,S+T,P,0,Z,C,V)$.
\end{theorem}

\begin{theorem}[Bounded register subtraction run]
\lean{Complexity.LogProg.TQBFARM.sb_run}
\label{Complexity.LogProg.TQBFARM.sb_run}
\leanok
\uses{Complexity.LogProg.TQBFARM.A, Complexity.LogProg.TQBFARM.Run, Complexity.LogProg.TQBFARM.St, Complexity.LogProg.TQBFARM.St.Bd, Complexity.LogProg.TQBFARM.St.Bd.toV, Complexity.LogProg.TQBFARM.St.toV, Complexity.LogProg.TQBFARM.orc, Complexity.LogProg.TQBFARM.rC, Complexity.LogProg.TQBFARM.rV, Complexity.LogProg.sub_run}
\difficulty{3}
Let $x$ be a finite bit string, let $B$ be a natural number, let $b$ be a Boolean value,
and let $st$ be a seven-register state whose registers $K,S,P,T,Z,C,V$ are all at most
$B$. The program has a finite run from label $\text{sb}(b,0)$ with registers initialized
from $st$ to label $\text{jzC}(b)$, with the same registers as $st$ except that $V$ is
set to $0$ and $C$ is set to the truncated difference $\text{max}(C-V,0)$; every
configuration before the endpoint satisfies the program guard and has all registers at
most $B$.
\end{theorem}

\begin{theorem}[Bounded halving routine]
\lean{Complexity.LogProg.TQBFARM.hv_run}
\label{Complexity.LogProg.TQBFARM.hv_run}
\leanok
\uses{Complexity.LogProg.TQBFARM.A, Complexity.LogProg.TQBFARM.Run, Complexity.LogProg.TQBFARM.St, Complexity.LogProg.TQBFARM.St.Bd, Complexity.LogProg.TQBFARM.St.Bd.toV, Complexity.LogProg.TQBFARM.St.toV, Complexity.LogProg.TQBFARM.orc, Complexity.LogProg.TQBFARM.rC, Complexity.LogProg.TQBFARM.rT, Complexity.LogProg.halve_run}
\difficulty{3}
Let $x$ be a finite bit string, $B$ a natural number, $b$ a Boolean, and $s$ a machine
state whose seven register values are at most $B$. The machine has a run on input $x$
from label $\text{hv}(b,0)$ and state $s$ to label $\text{od}(b)$ and state obtained from
$s$ by setting $C=0$ and $T=\left\lfloor s.T/2^{s.C}\right\rfloor$, with the run guard
satisfied at every configuration before its endpoint.
\end{theorem}
